%% cover_letter.tex — C14: Wide-Area GOOSE Consensus Protocol
%% IEEE Transactions on Smart Grid
\documentclass[11pt,a4paper]{letter}
\usepackage[T1]{fontenc}
\usepackage[utf8]{inputenc}
\usepackage[margin=2.5cm]{geometry}
\usepackage{microtype,amsmath}

\signature{Anoop~V.~Eluvathingal\\
           Ph.D.\ Candidate, Dept.\ of Electrical Engineering\\
           Indian Institute of Technology Madras\\
           Chennai 600036, India\\[2pt]
           \texttt{arun.eluvathingal@iitm.ac.in}}

\address{Department of Electrical Engineering\\
         Indian Institute of Technology Madras\\
         Chennai 600036, India}

\begin{document}
\begin{letter}{Editor-in-Chief\\
               \textit{IEEE Transactions on Smart Grid}}

\opening{Dear Editor,}

We submit for your consideration:

\begin{center}
  \textbf{Wide-Area Consensus via GOOSE for Coordinated Protection
  in IBR-Rich Distribution Networks}
\end{center}

\textbf{The problem.}
In IBR-penetrated distribution networks, a single substation may
see a bounded fault current below the local sensitivity threshold,
while adjacent substations observe clear fault signatures.  No
individual relay has the full picture.  Classical wide-area
protection schemes aggregate raw measurements but do not account
for the varying \emph{reliability} of estimates across substations
under IBR operating conditions.

\textbf{Contributions.}
\begin{enumerate}
  \item \textbf{Quality-weighted consensus (C14)}: Trip votes are
        weighted by $w_i = \max(0, 1 - k_n^{(i)}/30)$, derived
        directly from the SAMBPS condition number $k_n^{(i)}$.
        Ill-conditioned substations are down-weighted automatically,
        without any manual configuration.

  \item \textbf{Conservativeness preservation (Theorem 1)}: The
        multi-node consensus trip rule is proven to inherit the
        zero-false-trip guarantee of the single-node CUEP trip rule
        (C13), making the aggregate at least as conservative as any
        constituent relay.

  \item \textbf{Dropout robustness (Proposition 1)}: Zone accuracy
        $\geq 94.8\%$ with up to $\lfloor (N_{\rm SS}-1)/2 \rfloor$
        simultaneous substation dropouts, degrading gracefully
        to unweighted majority in the limit.

  \item \textbf{Hardware validation}: 6.1\,ms P99 trip latency on
        a three-relay IEC\,61850 Class P2 testbed, meeting the
        10\,ms requirement with 38\% margin.
\end{enumerate}

\textbf{Fit to \textit{IEEE Transactions on Smart Grid}.}
The Smart Grid journal covers coordinated control and communication
for power system protection and automation.  This paper directly
addresses the smart grid integration of IBR-aware protection
coordination, with a practical IEC\,61850 GOOSE implementation
and hardware demonstration.

\textbf{Suggested reviewers.}
\begin{itemize}
  \item Prof.\ Arun G.\ Phadke (Virginia Tech) --- wide-area protection
  \item Prof.\ Ali Hooshyar (University of Toronto) --- IBR protection
  \item Prof.\ Mladen Kezunovic (Texas A\&M) --- smart grid relaying
  \item Dr.\ Alexander Apostolov (OMICRON) --- IEC 61850, GOOSE
\end{itemize}

This manuscript is not under review elsewhere.
All authors have approved the submission.

\closing{Sincerely,}
\end{letter}
\end{document}
